% rn-accessibility.tex — React Native accessibility: what the props become natively, accessible /
% role / aria-* vs the older accessibility* props, grouping, focus order and moving focus, live regions
% (Android) vs announcements, user settings (reduce motion, bold text, screen reader), text scaling,
% TalkBack vs VoiceOver differences, testing (Accessibility Inspector, Accessibility Scanner,
% RNTL *ByRole), WCAG basics (contrast, target size).
% Sources (checked 2026-09-25):
%   https://reactnative.dev/docs/accessibility · https://reactnative.dev/docs/accessibilityinfo
%   https://reactnative.dev/blog/2026/02/11/react-native-0.84 (pressable Text gets role link)
%   https://reactnative.dev/blog/2026/04/07/react-native-0.85 (setAccessibilityFocus deprecated)
%   https://oss.callstack.com/react-native-testing-library/docs/api/queries (*ByRole rules)
%   https://www.w3.org/TR/WCAG22/ (1.4.3, 1.4.11, 2.5.8) · Apple HIG (44 pt) · Material (48 dp)
% Not repeated: VoiceOver anatomy (label/value/trait/hint), Dynamic Type in UIKit/SwiftUI
% (ios-platform/accessibility-localization).
% @labels: area=react-native kind=api platform=cross-platform topic=ui,testing,platform-apis discussion=196
% @tags: accessibility, aria-props, accessibilityrole, voiceover, talkback, live-region, accessibilityinfo, reduce-motion, font-scaling, touch-target, wcag, getbyrole
\documentclass[8pt]{extarticle}
\usepackage{printup-sheet}
\usepackage{array}

\lstdefinelanguage{TSXSheet}{
  morekeywords={const,let,return,function,true,false,useEffect,expect},
  morekeywords=[2]{Pressable,View,Text,AccessibilityInfo,screen,getByRole},
  keywordstyle=[2]{\color{sheetOrange}\bfseries},
  sensitive=true, morecomment=[l]{//}, morestring=[b]", morestring=[b]',
  literate={===}{{\hbox{=}\hbox{=}\hbox{=}}}3 {=>}{{\hbox{=}\hbox{>}}}2 {==}{{\hbox{=}\hbox{=}}}2
           {</}{{\hbox{<}\hbox{/}}}2 {/>}{{\hbox{/}\hbox{>}}}2}
\lstset{basicstyle=\ttfamily\scriptsize, aboveskip=2pt, belowskip=2pt}

\newcommand\ct[1]{\texttt{#1}}
\newcommand\hd[1]{\par\vspace{3pt}\noindent{\bfseries\color{sheetBlue}#1}\par\vspace{1pt}}

\tikzset{
  pr/.style={draw=sheetOrange, fill=sheetOrange!8, font=\tiny\ttfamily, minimum height=3.9mm,
             minimum width=27mm, inner sep=1pt, anchor=west},
  io/.style={draw=sheetBlue, fill=sheetBlue!6, font=\tiny, minimum height=3.9mm,
             minimum width=33mm, inner sep=1pt, anchor=west},
  an/.style={draw=sheetGreen!70!black, fill=sheetGreen!8, font=\tiny, minimum height=3.9mm,
             minimum width=38mm, inner sep=1pt, anchor=west},
  lbl/.style={font=\tiny, text=black!75, inner sep=1pt, align=center},
  el/.style={draw=sheetRed, thick, rounded corners=2pt, dashed},
}

\begin{document}

\sheettitle{React Native accessibility — props, platforms, testing}{react-native · folio}

\oneliner{VoiceOver and TalkBack read the \textbf{native} view tree, not your JSX: RN maps each
prop onto \ct{UIAccessibility} (iOS) or \ct{AccessibilityNodeInfo} (Android). Prefer the
web-aligned \ct{role} + \ct{aria-*} props (\ct{role} wins over \ct{accessibilityRole}); a few needs
stay platform-specific — \textbf{live regions} are Android-only, iOS uses \textbf{announcements};
modal scoping is iOS-only.}

\vspace{3pt}
\noindent\begin{tikzpicture}[sheet]
  % ---------- prop mapping ----------
  \node[font=\bfseries\small, text=sheetBlue, anchor=west] at (0,3.65) {What a prop becomes};
  \node[lbl, font=\tiny\bfseries, anchor=west] at (0,3.3) {RN prop};
  \node[lbl, font=\tiny\bfseries, anchor=west, text=sheetBlue] at (2.8,3.3) {iOS (VoiceOver)};
  \node[lbl, font=\tiny\bfseries, anchor=west, text=sheetGreen!50!black] at (6.2,3.3) {Android (TalkBack)};
  \foreach \i/\p/\a/\b in {
    0/{accessible}/{isAccessibilityElement}/{focusable},
    1/{role · accessibilityRole}/{traits (.button, .header\ldots)}/{node class / role description},
    2/{aria-label}/{accessibilityLabel}/{contentDescription},
    3/{aria-disabled/-checked/\ldots}/{traits .notEnabled, .selected}/{node state (enabled, checked\ldots)},
    4/{aria-live="polite"}/{\textcolor{sheetRed}{none — call announce}}/{live region},
    5/{aria-modal}/{accessibilityViewIsModal}/{\textcolor{sheetRed}{none — hide siblings}},
    6/{aria-hidden}/{accessibilityElementsHidden}/{importantForA11y no-hide-descendants}}
  {
    \node[pr] at (0,2.9-\i*0.43) {\p};
    \node[io] at (2.8,2.9-\i*0.43) {\a};
    \node[an] at (6.2,2.9-\i*0.43) {\b};
  }
  % ---------- grouping ----------
  \draw[sheetGrey!40] (10.45,3.75) -- (10.45,0.1);
  \node[font=\bfseries\small, text=sheetBlue, anchor=west] at (10.6,3.65) {Grouping a list row};
  \node[lbl, anchor=west] at (10.6,3.3) {\textbf{before}: 3 swipes, no context};
  \draw[sheetGrey, rounded corners=2pt] (10.6,2.05) rectangle (16.5,3.15);
  \node[font=\scriptsize\bfseries] (a) at (11.2,2.85) {Alice};
  \node[font=\tiny] (b) at (11.55,2.3) {2 new messages};
  \node[draw=sheetBlue, fill=sheetBlue!10, font=\tiny, rounded corners=1pt] (c) at (15.6,2.6) {Open};
  \foreach \n/\k in {a/1,b/2,c/3} \node[el, fit=(\n), inner sep=1pt, label={[lbl, text=sheetRed]right:\k}] {};
  \node[lbl, anchor=west] at (10.6,1.75) {\textbf{after}: \ct{accessible} on the row + an action};
  \draw[sheetGrey, rounded corners=2pt] (10.6,0.55) rectangle (16.5,1.55);
  \node[font=\scriptsize\bfseries] at (11.4,1.2) {Alice};
  \node[font=\tiny] at (11.75,0.8) {2 new messages};
  \node[draw=sheetBlue, fill=sheetBlue!10, font=\tiny, rounded corners=1pt] at (15.6,1.05) {Open};
  \draw[el, draw=sheetGreen!60!black, solid] (10.65,0.6) rectangle (16.45,1.5);
  \node[lbl, anchor=west, text=sheetGreen!50!black, align=left] at (10.6,0.28)
       {1 swipe: ``Alice, 2 new messages, button'' · \emph{Actions}: Archive};
\end{tikzpicture}

\begin{multicols}{2}
\raggedright\setstretch{1.0}

\hd{How it works}
\begin{itemize}
  \item \textbf{Elements by default}: \ct{Text}, \ct{TextInput}, \ct{Switch}; \ct{Pressable}/touchables
        render an \ct{accessible} View. A plain \ct{View} is not. An \ct{accessible} parent
        \textbf{groups} its children: one focus stop, label = the children's text joined.
  \item \ct{role} takes web names (\ct{heading}, \ct{img}, \ct{slider}, \ct{searchbox},
        \ct{list}, \ct{presentation}\ldots); state/value via \ct{aria-checked} (\ct{'mixed'}),
        \ct{aria-expanded}, \ct{aria-busy}, \ct{aria-valuemin/max/now/text} — same data as
        \ct{accessibilityState} / \ct{accessibilityValue}.
  \item \textbf{Hint}: iOS reads it after a pause and users can switch hints off; Android
        \textbf{always} reads it. Keep it a short result (``Removes it''), never the label again.
  \item \textbf{Actions}: \ct{accessibilityActions} (\ct{activate}, \ct{increment}/\ct{decrement}
        for adjustable, \ct{longpress}/\ct{expand}/\ct{collapse} Android, \ct{magicTap}/\ct{escape}
        iOS, or custom) + \ct{onAccessibilityAction} — the swipe-to-delete for screen readers.
  \item \textbf{Focus order} = native tree order. Reorder: \ct{experimental\_accessibilityOrder}
        (list of \ct{nativeID}s). Move focus after navigation or an error:
        \ct{AccessibilityInfo.sendAccessibilityEvent(ref, 'focus')} (\ct{setAccessibilityFocus}
        deprecated in 0.85).
  \item \textbf{Speak a change}: Android \ct{aria-live="polite"|"assertive"} on the view that
        changes; iOS has no live regions $\to$ \ct{announceForAccessibility(msg)} (iOS:
        \ct{\ldots WithOptions(msg, \{queue: true\})}, \ct{announcementFinished} event).
  \item \textbf{User settings} (Promises + change events): both \ct{isScreenReaderEnabled},
        \ct{isReduceMotionEnabled}; iOS \ct{isBoldTextEnabled}, \ct{isReduceTransparencyEnabled},
        \ct{isInvertColorsEnabled}; Android \ct{isHighTextContrastEnabled},
        \ct{getRecommendedTimeoutMillis} (toasts, auto-dismiss).
  \item \textbf{Text size}: \ct{allowFontScaling} (default \ct{true}), \ct{maxFontSizeMultiplier};
        read \ct{useWindowDimensions().fontScale}. Fixed-height text rows clip at 200\%.
\end{itemize}

\hd{VoiceOver vs TalkBack}
{\footnotesize
\begin{tabular}{@{}>{\raggedright\arraybackslash}p{15mm}>{\raggedright\arraybackslash}p{27mm}>{\raggedright\arraybackslash}p{28mm}@{}}
\toprule
& \textbf{VoiceOver} & \textbf{TalkBack} \\ \midrule
hint & optional, after pause & always read \\
live change & announce API & \ct{aria-live} region \\
modal scope & \ct{aria-modal} & hide siblings (\ct{aria-hidden}) \\
labelled by & — & \ct{aria-labelledby} (\ct{nativeID}) \\
extra gestures & escape (2-finger Z), magic tap & actions menu \\
\bottomrule
\end{tabular}}

\columnbreak

\hd{Example}
\begin{lstlisting}[language=TSXSheet]
<Pressable role="button" aria-label="Delete draft"
  accessibilityHint="Removes it from this phone"
  aria-disabled={saving} disabled={saving}
  hitSlop={10} onPress={remove}>          // 24pt icon, 44pt target
  <TrashIcon />
</Pressable>
<Pressable role="button" onPress={open}       // one stop per row
  accessibilityActions={[{ name: 'archive', label: 'Archive' }]}
  onAccessibilityAction={e =>
    e.nativeEvent.actionName === 'archive' && archive()}>
  <Text>Alice</Text><Text>2 new messages</Text>
</Pressable>
<Text aria-live="polite">{status}</Text>    // Android
useEffect(() => {                            // iOS (works on both)
  AccessibilityInfo.announceForAccessibility(status); }, [status]);
// test: only accessibility elements match, hidden ones excluded
expect(screen.getByRole('button', { name: 'Delete draft' }))
  .toBeDisabled();
\end{lstlisting}

\hd{WCAG basics}
\begin{itemize}
  \item \textbf{Contrast} (1.4.3): text $\geq$ \textbf{4.5:1}, large text (18 pt / 14 pt bold) $\geq$
        3:1; UI parts and focus rings $\geq$ 3:1 (1.4.11). Never colour alone.
  \item \textbf{Target size}: Apple \textbf{44$\times$44 pt}, Material \textbf{48$\times$48 dp};
        WCAG 2.2 AA minimum 24$\times$24 CSS px (2.5.8). \ct{hitSlop} grows the touch area
        without changing layout.
  \item Respect \textbf{reduce motion} (swap slides for fades; Reanimated \ct{useReducedMotion()}),
        don't auto-dismiss faster than the recommended timeout.
\end{itemize}

\hd{Testing}
\textbf{Manual}: VoiceOver + Xcode \textbf{Accessibility Inspector} (audit); TalkBack + Android
\textbf{Accessibility Scanner} (targets, contrast, labels). \textbf{Automated}: RNTL
\ct{*ByRole} / \ct{*ByLabelText} / \ct{*ByHintText} — a \ct{View} matches by role only if
\ct{accessible}. \ct{testID} is for E2E (iOS \ct{accessibilityIdentifier}) — never spoken.

\hd{Traps}
\begin{itemize}
  \trap{An \ct{accessible} card with a button inside: the button is no longer its own stop — expose it as an action.}
  \trap{``Disabled'' shown by opacity only: the screen reader still says it works — set
        \ct{disabled}/\ct{aria-disabled}.}
  \trap{\ct{aria-live} does nothing on iOS; \ct{aria-modal} nothing on Android.}
  \trap{Icon-only \ct{Pressable} without \ct{aria-label}: there is no name to read.}
\end{itemize}

\hd{Remember}
\textbf{``Role, name, state — then focus, announce, scale, 44/48.''}


\end{multicols}

\noindent{\footnotesize\color{sheetGrey}\textit{Related:} accessibility-localization (VoiceOver
anatomy, Dynamic Type) · rn-testing-typescript (RNTL) · rn-animations-gestures (reduce motion) ·
rn-navigation (focus after navigation)}

\end{document}
